% 5-anexa-c.tex starts here
% \raggedbottom for the appendix: the intro paragraph shares a page with no
% float (the ER diagram is [p], on its own page), so \flushbottom from the body
% would stretch that short paragraph's vertical glue into big gaps. Previously
% set by the now-removed LaTeX-tooling and PAN-manifest appendices.
\raggedbottom

\section{Modelul de date}\label{anexa:model-date}

Schema completă a bazei de date --- cele 20 de modele \gls{prisma} --- este
redată în \figref{fig:er-diagram} ca diagramă entitate-relație. Diagrama este
generată automat din \texttt{schema.prisma} (cu \textit{prisma-dbml-generator})
și randată cu \textit{dbdocs.io}, rămânând astfel sincronizată cu modelul de
date efectiv. Cheile primare sunt marcate cu simbolul cheii, cheile externe cu
simbolul de legătură, iar conectorii dintre tabele redau relațiile dintre
entități; tabelele de joncțiune (de exemplu \texttt{UserDietaryRestriction},
\texttt{RecipeCollectionEntry}) sunt așezate lângă entitățile pe care le leagă.

\begin{figure}[p]
  \centering
  % Image is the dbdocs PDF export cropped to its content with `pdfcrop` ---
  % the raw export carries a full-canvas white border that wastes ~1/3 of the
  % page. Re-crop after any re-export. keepaspectratio + the height cap keep
  % the diagram inside the text block.
  \makebox[\textwidth]{%
    \includegraphics[width=\textwidth,height=0.9\textheight,keepaspectratio]{kookio-er-diagram}}
  \caption[Diagrama entitate-relație a bazei de date]{Diagrama
    entitate-relație a schemei de date: cele 20 de modele și relațiile dintre
    ele, generate din \texttt{schema.prisma} și publicate pe
    \textit{dbdocs.io}. Versiunea interactivă, navigabilă:
    \url{https://dbdocs.io/dinea.eduard/kookio-db} (parolă de acces:
    \texttt{FMI\_2026}).}
  \label{fig:er-diagram}
\end{figure}

% 5-anexa-c.tex ends here
